\section{Introduction}
%%%%% Note from KGT: KEEP IN MIND WHAT AN INTRO SHOULD CONTAIN! %%%%%%%
% source: https://dominika-langenmayr.net/wp-content/uploads/2022/09/Session-5_VL.pdf
% 1. Motivation (1 paragraph) , i.e. the hook
% 2. Research question  (1 paragraph) --> place it in the literatzre
% 3. Main contribution (2-3 paragraphs) , lead with our findings then how it extends the literature
% 4. Methods (1-2 pagraphs)
% 5. Findings (2-3 paragraphs) (optional, mention robustness)
% 6. Roadmap

The economic and social integration of refugees remains a central policy challenge across advanced economies \citep{brell_labor_2020}. Governments implement a range of policies not only to foster integration but also to shape migration flows, manage fiscal resources, and sustain public support for hosting refugees.\footnote{We use the term refugee for an individual whose asylum application was approved \citep{hatton2020}. We use the term asylum-seeker specifically for a person before being granted protection.} Some policies restrict refugees’ choices—such as dispersal rules that allocate them across regions—while others provide incentives, for instance by changing benefit generosity to foster labor market participation. Yet integration outcomes also depend on factors beyond policymakers’ control, including the absorptive capacity of local labor markets. Ultimately, successful integration hinges on how refugees respond to these policy environments and labor market conditions. Understanding these responses, and the ways in which different policies interact, is therefore critical for designing effective refugee integration strategies \citep{adda2022}.

We study how refugees adjust their labor supply and relocation decisions once restrictions on mobility and employment are lifted. Specifically, we examine how local labor market tightness and welfare benefit levels—individually and in interaction—shape refugees’ choices. Austria provides a particularly well-suited setting for this analysis. During the asylum process, refugees are subject to a dispersal policy and face strict limits on employment and relocation. These restrictions are lifted immediately upon recognition, at which point refugees gain full access to the labor market and freedom of movement across states. Three features make Austria distinctive. First, the sharp shift from restricted to unrestricted choice allows us to combine quasi-experimental variation in initial placement with observation of behavior under full mobility. Second, because most refugees must vacate their assigned accommodation within four months, they are forced to make consequential labor market and residential decisions in a compressed time frame, resembling the incentives embedded in work-first policies \citep{arendt_labor_2022, arendt_trade-offs_2023}. Third, Austria’s federal welfare system generates large and time-varying cross-state differences in benefit generosity, enabling us to test whether refugees respond to financial incentives by relocating to more generous states.\footnote{For example, in 2015, a single refugee with subsidiary protection in Carinthia received only €200 in basic subsistence support, compared to more than three times that amount in Tyrol or Vienna.}

We find that many refugees relocate immediately once restrictions are lifted: roughly 30 percent of those initially placed outside Vienna move to the capital within four months, despite Vienna having a relatively weak labor market. Both local labor market tightness and welfare generosity at the assigned location shape this decision. Tighter labor markets—especially in low-barrier occupations—induce higher and more immediate labor force participation among men, which in turn increases the probability of remaining in the assigned state. While these employment effects are temporary, location choices are persistently affected.

Benefit generosity also influences outcomes, though more modestly. A decrease in monthly benefits increases the propensity to relocate and male employment temporarily and female employment more persistently. Importantly, these effects interact with labor market conditions: benefit reductions significantly increase employment only when local labor markets are tight, consistent with recent evidence from Denmark \citep{dustmann2024}. Taken together, our results suggest that mobility responses primarily reflect adaptation to immediate constraints rather than forward-looking optimization over longer-term labor market prospects. Consequently, the combination of dispersal policies with subsequent unrestricted mobility does not contribute to "greasing the wheels" of the labor market \citep{borjas2001}.

Our contribution is threefold. First, we provide causal evidence on how local labor market conditions at the time of labor market entry shape refugees’ employment and mobility decisions. Second, we examine how cross-state differences in welfare benefit generosity influence both relocation and labor supply, thereby speaking to the welfare magnet hypothesis. Third, we analyze how local conditions and the policies in place interact, showing that the effects of welfare policies on employment are amplified when local labor markets are tight. Together, these findings inform academic debates on immigrant location choice and labor market assimilation, while also providing timely evidence for policymakers designing refugee integration policies.

We now turn to a more detailed description of our setting, data, methodology and results. We draw on the Austrian asylum system, which combines quasi-random assignment of asylum-seekers across regions with sharp changes in mobility and labor market access once protection is granted. This institutional structure creates exogenous variation in the local conditions refugees face at the time of labor market entry. Welfare benefits, set at the federal-state level, varied substantially across states, over time, and by protection status. Between 2015 and 2017, several states implemented large benefit cuts, some of which were later overturned by courts, further increasing variation in generosity. 

Our analysis uses administrative data covering more than 40,000 refugees who arrived between 2011 and 2018, providing complete information on employment trajectories, social security contributions, and place of residence. Unlike related studies \citep{ferwerda2023, dustmann2024, azlor_local_2020}, we can exploit both spatial and temporal variation in local labor market tightness and welfare benefit levels, which allows us to separate their effects and better account for regional confounders.

The data also contain detailed records on job openings, which enable us to measure both the overall tightness of the local labor market and the demand for specific occupations. These features of the Austrian setting allow us to examine how refugees respond to both labor market opportunities and financial incentives at the time their choice set expands.

An interquartile increase in local labor market tightness—about 0.2 additional vacancies per unemployed person—raises the probability of employment twelve months after recognition by 4.2 percentage points and reduces welfare use. The effect fades after roughly 30 months. Labor market tightness also increases the probability of remaining in the assigned state and decreases relocation to Vienna, the primary destination for movers. By contrast, lower benefit levels have modest effects: a reduction of €100 increases employment after one year by 0.5 percentage points and monthly earnings by about €10, while decreasing the probability of remaining in the assigned state by one percentage point. These employment effects dissipate after 2.5 years, but location effects persist. Neither labor market tightness nor benefits affect the probability of leaving Austria altogether. Finally, consistent with \citet{dustmann2024}, benefit reductions both raise employment and reduce welfare take-up only when local labor markets are tight.

%Austria received asylum applications at a rate above the EU average.\footnote{During 2014--2018, the EU average was two first applications per 1,000 inhabitants. Austria received four applications per 1,000 \citep{eurostat_2023}.} 
%We use administrative data, which contains the full labor market trajectories, social security insurance spells and place of residence for over 40,000 refugees arriving in Austria from 2011 to 2018. Identification is based on the initial exogenous distribution and the restrictions during the asylum process. Variation in labor demand and benefit levels over time and space allows us to identify their effects separately but also study how they interact. Unlike related studies \citep{ferwerda2023, dustmann2024, azlor_local_2020}, we can use spatial and temporal variation in labor demand \emph{and} benefit levels, which allows us to better control for potential regional confounders.

%The data includes comprehensive and detailed information on job openings, which helps us assess labor demand by quantity and type of jobs in their assigned regions. Additionally, refugees encounter heterogeneous welfare benefits set by federal states, which underwent several reductions between 2015 and 2017, though the courts later overturned some of these cuts. Consequently, benefit levels differ markedly across states, over time, and according to the type of protection granted. 

%Our main findings are as follows. An interquartile increase in local labor demand, equal to 0.2 additional open positions per unemployed person, increases the employment probability 12 months after protection by about 4.2 pp and consequently reduces welfare use. The effect dissipates after 30 months. Effects on earnings are driven by higher employment and not by higher wages. %, similar to the one found in \citet{barsbai2024} for family migrants in the US.
%Higher labor demand increases the probability of remaining in the state of assignment and reduces the probability of moving to Vienna, the preferred destination of refugees who move within Austria. The effect of higher benefits is mirrored but small: 100 euro higher benefit levels reduce the employment probability 12 months after protection by 0.5 pp. and monthly earnings by about 10 euro. Consequently, changes in benefit levels result in large changes in disposable income. However, the effects on employment dissipate after about 2.5 years. At the same time, 100 euro higher benefit levels persistently increase the probability of remaining in the state by about one pp. Neither local labor demand nor benefit levels affect the propensity to leave Austria entirely.

%Similar to \cite{dustmann2024}, we find that the effects of labor demand and benefit levels interact in significant ways. The effect of a reduction in benefit levels on employment after one year is large when labor demand is high and negligible otherwise. Conversely, lowering benefits when labor demand is high reduces the take-up of benefits but has no effect if labor demand is low.

%Our fine-grained data allows us to study the timing of events in detail. The effect of high labor demand materializes within three months. 

We contribute to several strands of literature that are usually studied in isolation but are nonetheless closely interlinked.

% Location choice of immigrants - in general and with respect to labor market and welfare generosity

\textbf{Location choice of immigrants.} The question of how immigrants choose a location has received extensive attention in the literature, but several issues remain contentious, such as the extent to which differences in welfare benefits and labor demand influence this decision. An empirical regularity observed in many contexts is that immigrants tend to cluster in large cities \citep{monras2023}. Our context is no exception, with widespread relocation to Vienna as soon as this becomes possible. This phenomenon has been explained by higher place utility offered by cities, for example, through the existence of ethnic enclaves or amenities such as educational institutions \citep{damm2010}. Another potential explanation is that migrants choose high-cost–high-income cities, as they consume a smaller share of their income locally and thus benefit disproportionately from higher wages \citep{albert2022}. The labor market is unlikely to play an important role as a pull factor in our context. Vienna is one of the locations with the lowest labor demand, especially for low-skilled jobs; average wages for blue-collar workers are below the Austrian average (€30,509 vs. €35,564);\footnote{\url{https://www.statistik.at/fileadmin/publications/Allgemeiner_Einkommensbericht_2020.pdf}} and refugees in Vienna have some of the lowest employment rates in the country. As the labor market does not seem to draw refugees to Vienna, other explanations—such as networks, amenities, and benefit generosity—are likely to play a more important role.

Given the strong concentration of refugees in a single destination and the absence of significant benefit changes in Vienna, we do not seek to identify which factors attract them to particular locations (pull) but rather which characteristics of their assigned locations push them away. Such push models have been used to study immigrants’ location preferences, but for causal identification they require exogenous placement into initial locations \citep{damm2009}, which is the case in our setting.

\citet{borjas2001} argues that migrants, who are less attached to a particular location and thus face lower moving costs, contribute to \textit{greasing the wheels} of the economy by avoiding areas where their skills are not in demand. Migrants reacting to local conditions through outmigration or staying is part of this process. Our results suggest that refugees do respond to the availability of jobs in their initial regions. However, they do not seem to select the subsequent regions based on labor market prospects. %However, these responses appear to ease short-term constraints rather than reflect longer-term optimization of earnings potential.

We also contribute to the question of whether the generosity of welfare benefits affects immigrants’ location choices—known as the \textit{welfare magnet hypothesis} \citep{borjas1999}. This question has received considerable attention. \citet{agersnap_welfare_2020} study the same Danish reform analyzed in \citet{dustmann2024} and show that the benefit reduction lowered migrant inflows with an implied elasticity of 1.3.\footnote{\citet{adema2025} discusses several methodological concerns with this study and argues that the elasticity is about 0.28 for new arrivals.} At the subnational level, \citet{ferwerda2023} study Switzerland and examine how the generosity of benefits in refugees’ initial placement locations affects subsequent relocation. This push approach is very similar to ours. They find that higher benefits reduce relocation, but only to a minor extent—at most two pp., even for very large differences in benefits. Our estimates of roughly one pp. for a €100 change in benefits are considerably larger. Exploiting some of the same reforms that we analyze, \citet{huber_impact_2022} find in an event-study design that reductions in benefit levels trigger significant outmigration. 



%The question of how immigrants choose a location has received extensive attention in the literature but several questions remain contentious, such as to what extent differences in welfare benefits and labor demand influence this decision. An empirical regularity observed in many contexts is that immigrants tend to cluster in large cities \citep{monras2023}. Our context is no exception with the widespread relocation to Vienna as soon as this is possible. This phenomenon has been explained with higher place utility offered by cities, e.g., through the existence of ethnic enclaves but also amenities such as educational institutions \citep{damm2010}. Another potential explanation is that migrants choose high-cost-high-income cities as they consume a smaller share of their income in this location and thus benefit disproportionately from the higher wages \citep{albert2022}. The labor market is unlikely to play an important role as a pull factor in our context. Vienna is one of the locations with the lowest labor demand, especially for low-skilled jobs; average wages for blue-collar workers are below the Austrian average (30,509 vs. 35,564 euro)\footnote{\url{https://www.statistik.at/fileadmin/publications/Allgemeiner_Einkommensbericht_2020.pdf}}; and refugees in Vienna have some of the lowest employment rates in the country.

% (see Section XX) ADD DESCRIPTIVES HERE
% CITE OECD REPORT HERE

%Given the strong concentration of refugees in one destination within Austria, we do not seek to identify which factors attract them to a certain destination (pull) but rather which factors of their assigned location push them away. Such push models have been used to learn about immigrants' location preferences, but for causal identification, they require exogenous placement into initial locations \citep{damm2009}, which is the case in our setup.

%\cite{borjas2001} argues that migrants, who are less attached to a particular location and thus have lower moving costs, contribute to \textit{greasing the wheels} of the economy by not moving to areas where their skills are not in demand. Migrants reacting to local conditions with outmigration/staying is part of greasing the wheels. Our results suggest that migrants do react to the availability of jobs in their initial region. However, these appear to relax their short-term constraints  and are not the result of longer-term optimization of their earnings potential.

%We also contribute to the question of whether the generosity of welfare benefits affects immigrants' location choices---known as the \textit{welfare magnet hypothesis} \citep{borjas1999}. This question has received considerable attention. \cite{agersnap_welfare_2020} study the same Danish reform as studied in \cite{dustmann2024} and show that the benefit reduction reduced migrant inflows with an implied elasticity of 1.3. At the sub-national level, \cite{ferwerda2023} study for Switzerland, how the generosity of benefits in the place where refugees were initially placed affects their subsequent relocation. This push approach is very similar to ours. They find that higher benefits decrease the incidence of relocation, but only to a minor extent. At most, by two pp., even for very large differences in benefits. Our estimates of a roughly one pp. change for a 100 euro change in benefits is considerably larger. Exploiting some of the same reforms that we look at, \cite{huber_impact_2022} find in an event study approach that reductions in benefit levels trigger significant outmigration.

% ANDREAS: I DROPPED THIS SECTION FOR NOW TO SHORTEN THE INTRO
%\textbf{Dispersal policies.} Many countries have introduced dispersal policies that distribute refugees according to some administrative rules. Through dispersal policies, countries seek to overcome the strong concentration of refugees in certain locations, use resources such as housing more evenly, and increase the acceptance of hosting refugees in the native population. Evidence on the effects of dispersal policies on refugee integration is scarce but the available evidence suggests that refugees subject to a dispersal policy are less integrated into the labor market than refugees who could choose their location themselves \citep{edin2004, fasani_struggle_2022}. While we cannot directly identify the effect of a dispersal policy, the Austrian context is interesting as it combines exogenous assignment through a dispersal policy during the asylum process with no relocation restrictions after getting protection.\footnote{This differs, for example, from the Swiss \citep{bansak_improving_2018} or Danish dispersal \citep{azlor2020} policies where refugees are prohibited or at least severely restricted to remain in the original location for multiple years after being granted protection.} Thus, the policy helps us to overcome the often observed self-selection of immigrants with favourable characteristics into locations with favourable labour market conditions \citep{azlor2020}. But at the same time, we can still observe the refugees' choices after they received protection. 

% Effects of dispersal policies

% many countries have dispersal policies. whats special about austria is that there is a dispersal policy but after status is granted, migrants are free to settle in the entire country. even more, KONs need to find new housing 4 months after status is granted.
% the idea of dispersal policies is to use resources more evenly but also to increase acceptance of hosting refugees in the native population.
%Valentin: in optimum, forward looking migrants should select a location based on their expected integration chances. however, it's more likely that the base their decisions on immediate opportunities and constraints.
% Several countries have started considering dispersal policies that place refugees already in the beginning according to individual integration chances.
% Municipalities could consider asking for refugees that fit their labor demand. however, a given refugee needs a lot of time before he/she gets access to the labor market.
%Migrants incur costs if they are at a suboptimal location initially and need to move. Local knowledge is lost
%We see huge effects for "preferred occupations"

%refugees have to make very consequential decisions within a short period of time.

%In a study for Sweden, Edin et al. (2004) show that, relative to refugees who arrived before the introduction of the policy (and were therefore free to choose where to settle), dispersed refugees were less likely to be employed, had lower earnings and claimed more welfare benefits. These findings are confirmed by Åslund and Rooth (2007). In the context of the Danish dispersal policy, instead, Damm and Rosholm (2010) provide evidence that relocation from the originally assigned location had a large positive effect on the probability of finding employment. 
% Our study is related to three existing impact evaluations that exploit spatial dispersal policies on refugees in Scandinavia to identify the causal effect of local labour market conditions on immigrant labour market outcomes: Åslund and Rooth (2007) for Sweden, Damm and Rosholm (2010) for Denmark and Godøy (2017) for Norway. While our study provides quasi-experimental estimates of the effects of the initial employment rate in the assigned location as well as for a range of alternative measures of local labour market conditions on immigrant employment, Åslund and Rooth (2007) estimates the effects of the initial local unemployment rate on both immigrant earnings and employment. Godøy (2017) , instead, estimates the effects of the initial regional employment rate of non-OECD immigrants on immigrant earn- ings, whereas Damm and Rosholm (2010) exploits the first Danish Spa- tial Dispersal Policy on refugees in place from 1986–1998 to provide quasi-experimental evidence on local determinants on the hazard rate into first job. (Azlor et al. 2020)
% 
%Higher non-western immigrant share decreases the likelihood of moving out. higher co-national share has the same effect but not precisely estimated. Commuting opportunities have no effect.

% refer to our model as a push model (see Hangartner et al 2023). An alternative way of learning about immigrants’ location preferences is to estimate push factors, i.e. the set of negative or positive social or economic factors in an area of origin which pushes migrants away (Lee 1966), based on immigrants’ subsequent internal migration pattern. However, push factor estimates may be biased if immigrants sort into initial locations based on unobserved personal attributes that also influence the secondary migration probability. (Piil Damm 2009)

%Four theories of migrant location choice: welfare magnet (Borjas), greasing the wheels (Borjas), high-income-high-cost (Monras), ethnic enclaves
%Thomas Bauer also has several papers on migrants' location choice.
%Vienna has the lowest labor demand and relatedly low wages, especially for blue-collar workers. In 2017, a blue-collar worker in Vienna had a gross income of 30.509 Euro, while the Austrian average was 35.564 Euro. 
%Vienna has lots of social housing but to qualify, refugees need to have lived in Vienna for more than two years -> no immediate solution.
%Introduce concept of place utility (Piil Damm and Rosholm (2010)). Location choice depends on place utility and earnings (income + benefits).


% highlight Dustman et al (2023), Hangartner et al (2023), and Fasani & Minale? as closest papers
% our mobility effect is somewhat larger than from hangartners push model. they find that mobility decreases by at most 2pp. when going from a very low to very high benefits.
% what makes our empirics better than hangartner, dustmann but also early papers on initial conditions is that we have both, variation in benefits and economic conditions over time. they only have it one dimension, requiring stronger assumptions for identification
% refugee migration is an urban phenomenon. Monras argues that migrants go to expensive cities as the consume a smaller part of their income locally. however, since the labor market in Vienna is so difficult for migrants, the argument does not really work in this context.



%Conditions are very similar to Denmark with refugees not allowed to work. But "The new legislation further strengthened the effort to allocate refugees across all municipalities in the country and made eligibility to means-tested welfare benefits during the 3-year introduction programme conditional on residing in the assigned municipality." 
% The most extensively studied context concerning the effects of dispersal policies and refugee integration more generally is Denmark (CITE XXX). However, the Danish and Austrian context differ in a key dimension. In Austria, asylum-seekers are placed in a particular municipality and are free to settle in the entire country upon receiving protection status. In Denmark, asylum-seekers are hosted in remote reception centers during the asylum process and are only placed in a municipality once they receive protection status. They are then required to remain in the municipality for at least three years, otherwise they loose their protection status. Asylum-seekers can express preferences over locations but these can only be considered when municipality quotas are not yet binding (Azlor et al. 2020).
%Note that the first Danish resettlement policy was different and more similar to the Austrian context without linking residence in the assigned municipality and receipt of benefit. But refugees were assisted in finding housing. Piil Damm and Rosholm (2010) exploit this dispersal policy.

% Effects of initial conditions

\textbf{Effect of initial labor market conditions.} Our findings on the effect of labor market conditions when entering the labor force complement active literature on initial conditions' (long-term) effects for refugees \citep{aslund_when_2007, godoy_local_2017, fasani_struggle_2022, aksoy_first_2023} and immigrants more generally \citep{chiswick1997, chiswick2002, barsbai2024}.\footnote{\citet{von_wachter_persistent_2020} provides a summary of the findings of graduating during different points of the business cycle on life outcomes: On average, entering the labor market in a recession reduces earnings, and these effects persist for up to 15 years, with large heterogeneities by socio-demographics (e.g., education, race, sex). For Austria, \citet{brunner_impact_2014} find that for labor market entrants during 1978-2000, a one pp. higher initial local unemployment rate decreases wages by about 1.3\% over the life cycle and persists for 20 years after labor market entry.} Our detailed administrative data on vacancies and unemployment allows for several refinements compared to previous work. First, we consider the conditions in the district and time of receiving unrestricted labor market access instead of at arrival and thus better capture the relevant moment.\footnote{\citet{aslund_when_2007} discuss which is the relevant initial labor market area since refugees do not immediately enter the labor market.} Second, we use the initial vacancy-to-unemployment ratio ($IVUR$) instead of employment or unemployment ratios because it might be a better measure of local labor market conditions \citep{dustmann2024}. In particular, the measure provides us not only with information on the general state of the labor market but we can identify labor demand for particular skills or individual sectors. Third, variation over time and space allows us to not only rely on cross-sectional variation but also to control for other unobserved regional characteristics that could be correlated with $IVUR$ and confound our results. Fourth, we can also look at the effects on a monthly level, finely capturing any dynamics playing out immediately after getting access to the labor market. Two concurrent papers \citep{degenhardt2025initial, degenhardt2025} exploit the same Austrian dispersal setting to study single labor-market segments (hospitality and gig work, respectively); we compare their findings with ours in Section \ref{sec:mech} alongside our own sectoral decompositions.
Relative to the majority of findings in this literature, we find that the effects of initial labor market conditions on employment are not persistent beyond 2.5 years. This is likely due to a situation where labor supply decisions and location choice are made simultaneously. Conditions right after receiving protection are important to find employment and subsequently housing.\footnote{The only other study we are aware of to analyze the effects on mobility is \citet{barsbai2024}, which finds that family migrants in the U.S. do not move in response to adverse economic conditions. However, unlike refugees, family migrants have strong personal links to the location of initial residence, making any moves more costly.}
%Refugees who face strong labor demand 
Several factors might explain our short-lived effects. First, the good labor market conditions at arrival could push refugees into employment too fast and `hinder' investments in host-country-specific human capital, which would, otherwise, payout in the future \citep{battisti_dynamic_2022}. Second, refugees in Austria are particularly mobile, and the importance of initial local conditions fades away. We also study how initial labor market tightness affects mobility and find persistent effects on the probability of staying in the assigned location. 


%Next, more specifically, we contribute to the literature studying the role of initial local (economic) conditions on refugees` integration outcomes. These initial conditions cover a variety of items and can be grouped as follows \citep{muller_labor_2023}: labor market conditions (e.g., unemployment rates \citep{aslund_when_2007, godoy_local_2017, aksoy_first_2023, muller_labor_2023}, recessions \citep{barsbai2024}), co-ethnic networks \citep{beaman_social_2012, edin_ethnic_2003, battisti_dynamic_2022}, and asylum policies (e.g., employment bans \citep{fasani_lift_2021}, dispersal policies \citep{fasani_struggle_2022}), and local (negative) attitudes \citep{aksoy_first_2023, jaschke_scared_2022, muller_labor_2023}.
%From these papers, those studying the role of unemployment rates have usually overcome selection problems by relying on quasi-random variation provided by, e.g., dispersal policies \citep{aslund_when_2007, godoy_local_2017, aksoy_first_2023, muller_labor_2023} or family reunification policies \citep{barsbai2024}.
%For Sweden, \citet{aslund_when_2007} study the long-term effects of initial local unemployment rates on refugees arriving between 1987 and 1991. For identification, they rely on a refugee settlement policy and find that initial higher local unemployment rates have a long-lasting negative impact on employment and earnings. These appear to be a combination of scarring and lock-in effects. The authors do not only use the initial local unemployment rate but also the one in the second full year after arrival since refugees do not immediately enter the labor market, where the latter is larger than the former.
%\citet{godoy_local_2017} study the role of local conditions---measured by the employment rate of the immigrant population in the initial labor market region---at the time of arrival on labor earnings of refugee immigrants in Norway. She relies on the quasi-random allocation of refugees to identify the effect of the initial immigrant unemployment rate. Interestingly, she finds that the effects are not driven by individual scarring effects but rather by persistent local poor economic conditions together with immobility.
%\citet{muller_labor_2023} study the role of initial conditions (unemployment rate, co-ethnic network, restrictive attitudes) on the integration of refugees in Switzerland during 1998-2018. They study the interplay of local contexts and policies. The authors find significant negative effects of initial unemployment rates and positive effects of initial more restrictive attitudes toward refugees and immigrants on the employment probabilities of refugees. These effects persist for over 16 years after arrival in Switzerland.
%\citet{aksoy_first_2023} study the role of unemployment rates (fixed in 2014) and attitudes towards immigrants on the labor market and social integration of refugees who arrived in Germany during 2013-2018. Using the German socio-economic panel they find that both the unemployment rate and negative attitudes decrease the employment probability of refugees. For identification, the authors rely on the quasi-random allocation of refugees across counties. %, which they validate by showing no correlation between individual-level characteristics and the county-level unemployment rate and attitudes toward immigration.
%Our paper contributes to this literature in several ways. First, we rely not only on the dispersal policy in place but acknowledge that municipalities might prefer to host certain groups. Hence, conditional on arrival-group fixed-effects% (family status*gender*quarter of arrival)


% Effect of welfare generosity


\textbf{Effect of welfare generosity.} Moreover, our work also contributes to a small literature studying the link of welfare generosity and labor market conditions on the integration and welfare use by immigrants \citep{Bratsberg2020ImmigrantGenerosity}. In the most closely related paper, \citet{dustmann2024} study the role of refugee benefit cuts in Denmark in an RD setting. The initial positive employment effects disappeared after 4-5 years. The results further show that the effect is larger in regions with more favorable labor demand conditions. Similar to this study, we also identify the interaction of the effect of these two factors. Our findings corroborate the conclusion of \citet{dustmann2024} as we also see that changes in benefit levels have larger effects on employment rates and welfare use when labor markets are tight.

More broadly, our paper expands on previous literature by looking at the interaction of local conditions and the policies in place. Although several studies have recognized the need for this approach \citep[e.g.,][]{fasani_struggle_2022, muller_labor_2023}, few have incorporated both elements in their analysis. Recent work has examined the combined role of initial labor market conditions and hostile attitudes toward immigrants \citep{muller_labor_2023, aksoy_first_2023, schilling_impact_2021}, but aside from \citet{dustmann2024}, none have focused on the interaction with welfare policies.






%XXXXXXXXXXX




%Since 2015, Europe and, more recently also, the U.S. have seen a sharp increase in asylum applications (UNHCR, 2022). Labor market integration of refugees is of particular importance, with employment rates of refugees often significantly lower than those of natives and other migrant groups \citep{brell2020}. This may result in extended periods of reliance on welfare benefits. Recently, many countries have restricted access to or reduced welfare benefits for refugees. This strategy aims to incentivize greater labor force participation but also to make these destinations less appealing. However, labor demand also plays a crucial role, potentially affecting both labor market integration and the attractiveness of a destination.



% SOME NOTES (ANDREAS)



%Still, the big question is why the "headstart" migrants get in some locations is not converted into a longer-term advantage.
%Danish policy seeks to integrate refugees into low-skilled jobs (Foged and Peri). This might be counterproductive in the longer term.
%Our high-frequency data allows us to study the immediate reaction. The effect of high labor demand occurs almost immediately - after 3 months, the coefficient is already around 0.15. Which is huge relative to the very low baseline. Also the reaction to benefit levels is quick and happens within three months. As a consequence, they stay - we see the immediate location choice influenced.
%even though we see that employment rates are not permanently shifted, location choice is. why? in this initial period migrants need to find an apartment. when they find a job quickly, they also solve the housing question. and then, even if the employment is not permanent, they have settled.
%add scatterplot between IVUR and staying rates


%FURTHER LITERATURE
%For investigations into the effect of ethnic concentration on immigrant labour market outcomes, see Borjas (1995), Beaman (2012), Cutler et al. (2008), Dagnelie et al. (2017), Damm (2009) and Edin et al. (2004), among others.  (from the Minale/Fasani paper)

%Effect of employment rate at arrival is constant for four years.
%The findings of a significant, negative effect of the presence of immigrants and the local population size on the hazard rate into first job both support the assumption underlying dispersal policies, that spatial dispersal of refugees away from immigrant-dense cities promotes economic assimilation of refugees. If the aim of a spatial dispersal policy is to pursue the assimilation motive—the policy should primarily be designed so as to reduce relocation rates while at the same time placing refugees in municipalities where the job finding rate is fairly high. Medium-sized municipalities with access to institutions for vocational and higher education, an appropriate number of co-nationals (but no immigrants of other ethnicities), a large public housing market and a low unemployment rate appear to be well suited for this purpose. (Piil Damm and Rosholm (2009)
%To summarize, refugees prefer living in large cities in part because it facilitates access to housing and educational institutions. (Piil Damm 2009)



%Although we cannot directly test for the quality of the first job, the effect on the log real monthly wage income (conditional on employment) is positive---although, statistically not significant---only during the first 12 months after protection. The effects then fade away, suggesting low-quality initial jobs.

%We find no evidence of either $IVUR$ or $IPBL$ affecting the use of services from the public employment service. A higher $IVUR$ decreases the probability of receiving welfare benefits. This effect almost mirrors the employment effects, indicating that a refugee is less likely to rely on welfare and profit from the favorable labor demand conditions by finding a job. We find no effect of higher $IPBL$ on the likelihood of receiving welfare benefits.

%we study the effects of labor demand and welfare benefit level (IPBL) at the time of labor market entry. We measure initial labor demand by the vacancy-to-unemployment ratio (IVUR) at the district of residence at the time an asylum-seeker got her asylum claim approved. Benefit levels vary substantially between federal states and by protection status. Our empirical strategy builds on the dispersal policy in place that distributes asylum-seekers according to quotas. Municipalities can express preferences about demographic characteristics of asylum-seekers (family status and gender). Conditional on arrival-group fixed-effects (arriving alone*gender*quarter of arrival*initial state) the allocation of asylum-seekers occurs in a haphazard way. Our identifying variation stems from cross-sectional within-group variation, i.e., some refugees with similar characteristics are assigned to districts with higher/lower $IVUR$ and $IPBL.$  Since asylum-seekers cannot influence the duration of their asylum process, we also use longitudinal variation across districts-of-protection. Hence, we can condition on district-of-protection FE to control for unobserved time-constant characteristics of the district that might be correlated with $IVUR$ and $IPBL.$  

%

%Among the major obstacles for labor market integration  are long asylum processes, lack of certificates and language proficiency, and the policies in place \citep[e.g.,][]{brell_labor_2020, becker_consequences_2019}. Many factors have been studied in isolation and in different settings, restricting the possibility for comparing their importance and for studying how different factors interact. 

%For example, the dispersal policies already determine local initial labor market conditions, access to networks, and to different levels of welfare benefits. While the role of initial labor market conditions, training, language courses and networks have been broadly studied in the literature, less is known about the role of policies in place, such as access to welfare benefits and its interaction with labor demand .



  %To address the selection of labor market entrants, they employ a standard instrument in this literature, which is to use an earlier local unemployment rate (e.g., at age 16) for the unemployment rate at the time of labor market entry. 
%We use the same core dataset (the ASSD) but focus on a different period (2011-2021) and a different sample: We study refugees instead of native male labor market entrants, and use a different empirical strategy.

%, we assume that the allocation of refugees was as good as random. %Finally, we additionally look at the role of benefit levels in place at the state and year of labor market access. This is novel to the literature since we can explore the role of potential (de)incentives to supply labor.

%We also contribute to a larger literature on immigrant labor market assimilation. \cite{Bratsberg2014ImmigrantsInsurance} compare different migrant groups and find that refugees assimilated quickly in early years but had rising social insurance rates and falling employment rates in the long run in Norway. %TO BE DISCUSSED HERE.

The rest of the paper is organized as follows. Section \ref{sec:background} provides information on the Austrian asylum system, the dispersal policy, and the welfare benefits for refugees. Section \ref{sec:dataoob} describes the data, sample, and key variables. Section \ref{sec:empirics} describes the empirical strategy. Section \ref{sec:results} presents the results, and Section \ref{sec:conclusion} concludes.%\footnote{This is a preliminary draft and not all sections are fully developed. We apologize for any ambiguities resulting from this.}


%And third, the assimilation of immigrants and refugees… \textcolor{red}{I still need to write this paragraph :)}
